\documentclass[10pt]{extarticle}
\usepackage[a4paper,margin=1.5cm]{geometry}
\usepackage[T1]{fontenc}
\usepackage{amsmath,amssymb}
\usepackage{graphicx}
\usepackage[export]{adjustbox}
\usepackage{multicol}
\usepackage{enumitem}
\usepackage{booktabs,array,tabularx}
\setlength{\columnsep}{0.7cm}
\setlength{\parindent}{0pt}
\setlength{\parskip}{3pt}
\pagestyle{empty}
\makeatletter
% figure registry (figdefs.tex is generated by the extract step)
\newcommand{\qfigdef}[4]{\@namedef{qf@file@#1}{#2}\@namedef{qf@w@#1}{#3}\@namedef{qf@wide@#1}{#4}}
\newcommand{\qf@check}[1]{\@ifundefined{qf@file@#1}{\PackageError{score2}{Unknown figure id #1}{}}{}}
\newcommand{\qf@img}[4][]{\includegraphics[width=#2\dimexpr\@nameuse{qf@w@#4}pt\relax,max width=#3,max height=0.45\textheight,#1]{\@nameuse{qf@file@#4}}}
% Full-width questions: finalize wraps every question that contains a wide figure in \qfullbegin...\qfullend,
% so the WHOLE question (stem, figures, options) is set across the page in reading order, then the two columns
% resume. The question's head - its number, stem and first figure (\qfullsplit marks the end of the head) -
% is kept on one page, so a number is never left at a page foot with its figure overleaf; the rest may break
% across pages normally, so no big gaps are left. Without a \qfullsplit, the whole question is kept together
% if it is under 45% of a page. Never write these yourself.
\newif\ifqf@full
\newif\ifqf@boxed
\newsavebox{\qf@box}
\newcommand{\qf@open}{\global\qf@fulltrue\par\global\qf@boxedtrue
  \setbox\qf@box\vbox\bgroup\hsize\textwidth\linewidth\textwidth}
\newcommand{\qf@emit}[1]{\par\egroup\global\qf@boxedfalse
  \ifdim\dimexpr\ht\qf@box+\dp\qf@box\relax<#1\textheight\noindent\usebox{\qf@box}\par\else\unvbox\qf@box\fi}
\newcommand{\qfullsplit}{\ifqf@boxed\qf@emit{0.85}\fi}
\newcommand{\qf@close}{\ifqf@boxed\qf@emit{0.45}\fi\global\qf@fullfalse}
\newcommand{\qfullbegin}{\end{multicols}\qf@open}
\newcommand{\qfullend}{\qf@close\begin{multicols}{2}}
\newcommand{\qfullstart}{\qf@open}
\newcommand{\qfulldone}{\qf@close}
% \qfig[scale]{id}: block figure, centred, directly where it is written (size set by the tool).
% Fallback if a wide figure is not inside a full-width question: break out of the columns in place - no floats.
\newcommand{\qfig}[2][1]{\qf@check{#2}%
  \ifqf@full\qf@blk{#1}{#2}\else
  \ifnum\@nameuse{qf@wide@#2}=1 \expandafter\@firstoftwo\else\expandafter\@secondoftwo\fi
  {\end{multicols}\begin{center}\qf@img{#1}{\textwidth}{#2}\end{center}\begin{multicols}{2}}%
  {\qf@blk{#1}{#2}}\fi}
\newcommand{\qf@blk}[2]{\par\smallskip{\centering\qf@img{#1}{\linewidth}{#2}\par}\smallskip}
% \qfigi[scale]{id}: inline figure (inside option lists / tables / minipages); never breaks columns.
\newcommand{\qfigi}[2][1]{\qf@check{#2}\qf@img[valign=c]{#1}{\linewidth}{#2}}
\makeatother
% structure macros - lint reads these to check completeness and figure placement
\newcommand{\papersec}[1]{\par\bigskip{\centering\Large\bfseries PAPER-#1\par}\medskip}
\newcommand{\partsec}[2]{\par\medskip\noindent{\large\bfseries PART-#1: #2}\par\smallskip}
\newcommand{\sect}[1]{\par\smallskip\noindent\textbf{#1}\par}
\newcommand{\qn}[1]{\par\smallskip\noindent\textbf{#1.}~}
\newcommand{\paraq}[2]{\par\smallskip\noindent\textbf{Paragraph for Questions #1 and #2}\par}
\newenvironment{qoptions}{\begin{itemize}[leftmargin=1.8em,itemsep=1pt,topsep=2pt,label={}]}{\end{itemize}}
\newcommand{\opt}[2]{\item[(#1)] #2}
% figure registry (generated by score2_tools.py finalize)
\qfigdef{p1_phy_s0_q01_1}{SRGBT - Test 03 (08 Aug 2021)-img/p1_phy_s0_q01_1.png}{114.0}{0}
\qfigdef{p1_phy_s0_q02_1}{SRGBT - Test 03 (08 Aug 2021)-img/p1_phy_s0_q02_1.png}{156.7}{0}
\qfigdef{p1_phy_s0_q03_1}{SRGBT - Test 03 (08 Aug 2021)-img/p1_phy_s0_q03_1.png}{71.2}{0}
\qfigdef{p1_phy_s0_q04_1}{SRGBT - Test 03 (08 Aug 2021)-img/p1_phy_s0_q04_1.png}{43.5}{0}
\qfigdef{p1_phy_s0_q05_1}{SRGBT - Test 03 (08 Aug 2021)-img/p1_phy_s0_q05_1.png}{88.6}{0}
\qfigdef{p1_phy_s0_q06_1}{SRGBT - Test 03 (08 Aug 2021)-img/p1_phy_s0_q06_1.png}{107.5}{0}
\qfigdef{p1_phy_s0_q07_1}{SRGBT - Test 03 (08 Aug 2021)-img/p1_phy_s0_q07_1.png}{126.9}{0}
\qfigdef{p1_phy_s0_q08_1}{SRGBT - Test 03 (08 Aug 2021)-img/p1_phy_s0_q08_1.png}{90.4}{0}
\qfigdef{p1_phy_s0_q09_1}{SRGBT - Test 03 (08 Aug 2021)-img/p1_phy_s0_q09_1.png}{81.4}{0}
\qfigdef{p1_phy_s0_q10_1}{SRGBT - Test 03 (08 Aug 2021)-img/p1_phy_s0_q10_1.png}{56.9}{0}
\qfigdef{p1_phy_s0_q11_1}{SRGBT - Test 03 (08 Aug 2021)-img/p1_phy_s0_q11_1.png}{68.4}{0}
\qfigdef{p1_phy_s0_q13_1}{SRGBT - Test 03 (08 Aug 2021)-img/p1_phy_s0_q13_1.png}{102.9}{0}
\qfigdef{p1_phy_s0_q14_1}{SRGBT - Test 03 (08 Aug 2021)-img/p1_phy_s0_q14_1.png}{121.0}{0}
\qfigdef{p1_phy_s0_q15_1}{SRGBT - Test 03 (08 Aug 2021)-img/p1_phy_s0_q15_1.png}{145.3}{0}
\qfigdef{p1_phy_s0_q16_1}{SRGBT - Test 03 (08 Aug 2021)-img/p1_phy_s0_q16_1.png}{85.2}{0}
\qfigdef{p1_phy_s0_q18_1}{SRGBT - Test 03 (08 Aug 2021)-img/p1_phy_s0_q18_1.png}{149.3}{0}
\qfigdef{p1_phy_s0_q19_1}{SRGBT - Test 03 (08 Aug 2021)-img/p1_phy_s0_q19_1.png}{104.8}{0}
\qfigdef{p1_phy_s0_q20_1}{SRGBT - Test 03 (08 Aug 2021)-img/p1_phy_s0_q20_1.png}{127.8}{0}
\qfigdef{p1_chem_s0_q21_1}{SRGBT - Test 03 (08 Aug 2021)-img/p1_chem_s0_q21_1.png}{294.8}{1}
\qfigdef{p1_chem_s0_q22_1}{SRGBT - Test 03 (08 Aug 2021)-img/p1_chem_s0_q22_1.png}{271.9}{1}
\qfigdef{p1_chem_s0_q24_1}{SRGBT - Test 03 (08 Aug 2021)-img/p1_chem_s0_q24_1.png}{160.7}{0}
\qfigdef{p1_chem_s0_q27_1}{SRGBT - Test 03 (08 Aug 2021)-img/p1_chem_s0_q27_1.png}{102.0}{0}
\qfigdef{p1_chem_s0_q37_1}{SRGBT - Test 03 (08 Aug 2021)-img/p1_chem_s0_q37_1.png}{281.5}{1}
\qfigdef{p1_chem_s0_q39_1}{SRGBT - Test 03 (08 Aug 2021)-img/p1_chem_s0_q39_1.png}{66.6}{0}
\raggedcolumns
\makeatletter\@beginparpenalty=10000\makeatother
\newcolumntype{L}[1]{>{\raggedright\arraybackslash}p{#1}}
\begin{document}
\begin{center}{\Large\bfseries SRGBT --- Test 03 (08 Aug 2021)}\end{center}
\begin{multicols}{2}
\raggedcolumns
\papersec{1}
\partsec{1}{PHYSICS}
\qn{1} A plank is moving along a smooth surface with a constant speed V. A block of mass M is gently placed on it. Initially the block slips and then acquires the constant speed (V) same as the plank. Throughout the period, a horizontal force is applied on the plank to keep its speed constant.\par
\qfig{p1_phy_s0_q01_1}
\begin{qoptions}
\opt{A}{the work performed by the external force is MV$^2$}
\opt{B}{the work performed by the external force is 1/2 MV$^2$}
\opt{C}{the heat developed due to friction between the block and the plank is 1/2 MV$^2$}
\opt{D}{the heat developed due to friction between the block and the plank is 1/4 MV$^2$}
\end{qoptions}
\qn{2} A 4m long rod AB slides down the plane with v$_\mathrm A$ = 4 m/s to the left and a$_\mathrm A$ = 5 m/s$^2$ to the right. The value of $\theta$ at given instant is 30$^\circ$. [cos 15$^\circ$ = 0.965; sin 15$^\circ$ = 0.258]\par
\qfig{p1_phy_s0_q02_1}
Choose the correct statement(s)
\begin{qoptions}
\opt{A}{The angular velocity $\omega$ of the rod at given instant of time is 0.73 rad/sec.}
\opt{B}{The angular velocity $\omega$ of the rod at given instant of time is 0.23 rad/sec.}
\opt{C}{The angular acceleration $\alpha$ of the rod at given instant of time is 0.234 rad/sec$^2$.}
\opt{D}{The angular acceleration $\alpha$ of the rod at given instant of time is 0.77 rad/sec$^2$.}
\end{qoptions}
\qn{3} Figure shows a smooth cylindrical pulley of radius R with centre at origin of co-ordinates. An ideal thread is thrown over it on the two parts of ideal thread two identical masses are tied initially at rest with co-ordinates (R, 0) and ($-$R, $-$R) respectively. If mass at x-axis is given a slight upward jerk, it leaves contact with pulley at (R cos$\phi$, Rsin$\phi$). Then ($\phi$ is the angular position of A)\par
\qfig{p1_phy_s0_q03_1}
\begin{qoptions}
\opt{A}{$\dfrac{\phi}{\sin\phi}=1.5$}
\opt{B}{$\dfrac{\phi}{\sin\phi}=2$}
\opt{C}{After rotating by angle 30$^\circ$ by particle A on the pulley speed of the particles is $\sqrt{\mathrm{gR}\left(\dfrac\pi3-1\right)}$}
\opt{D}{After rotating by angle 30$^\circ$ by particle A on the pulley speed of the particles is $\sqrt{\mathrm{gR}\left(\dfrac\pi6-\dfrac12\right)}$}
\end{qoptions}
\qn{4} A light cord of length $\ell$ is attached to a nail driven into curved surface of a cylinder A of radius r($\ell$ $>$ $\pi$r). The cylinder is fixed with its axis horizontal. A small ball B of mass m is hanging at the lower end of the cord. If V is horizontal velocity imparted to the ball B so that the cord will certainly slack during the subsequent motion is\par
\qfig{p1_phy_s0_q04_1}
\begin{qoptions}
\opt{A}{V $>\sqrt{2\mathrm g\ell}$}
\opt{B}{V $>\sqrt{2\mathrm g(\ell-\mathrm r)}$}
\opt{C}{V $<\sqrt{\mathrm g(5\ell-3\pi\mathrm r)}$}
\opt{D}{V $<\sqrt{5\mathrm g(\ell-4\pi\mathrm r)}$}
\end{qoptions}
\qn{5} Two particles of mass m each are kept on a horizontal circular platform on two mutually perpendicular radii at equal distance r from the center of the table. The particles are connected with a string, which is just taught when the platform is not rotating. The coefficient of static friction between the platform and block is $\mu$ (now if angular speed of platform m is slowly increased).\par
(If $\mu=\dfrac1{\sqrt2}$, r = 2.5 m, m = $\sqrt2$ kg and g = 10 m/s$^2$)\par
\qfig{p1_phy_s0_q05_1}
\begin{qoptions}
\opt{A}{The maximum angular speed ($\omega$) of platform about its center so that the blocks remain stationary relative to platform is 1 rad/sec.}
\opt{B}{The maximum angular speed ($\omega$) of platform about its center so that the blocks remain stationary relative to platform is 2 rad/sec.}
\opt{C}{Tension in the string when angular speed is maximum (above case) is 5 N.}
\opt{D}{Tension in the string when angular speed is maximum (above case) is 10 N.}
\end{qoptions}
\qn{6} A circular tube of mass M is placed vertically on a horizontal surface as shown in the figure. Two small spheres, each of mass m, just fit in the tube are released from the top, as shown in the figure. If $\theta$ gives the angle between radius vector of either ball with the verticall, if the tube breaks its contact with ground when $\theta$ = 60$^\circ$ (Ignore any friction) then\par
\qfig{p1_phy_s0_q06_1}
\begin{qoptions}
\opt{A}{ratio $\dfrac{\mathrm m}{\mathrm M}$ = 2}
\opt{B}{ratio $\dfrac{\mathrm m}{\mathrm M}$ = 3}
\opt{C}{relative acceleration of 1 particle w.r.t. another particle at angle $\theta=\cos^{-1}\dfrac23$ is $\dfrac{\sqrt3\mathrm g}2\left(1-\dfrac{\sqrt3}2\right)$ m/s$^2$.}
\opt{D}{relative acceleration of 1 particle w.r.t. another particle is $\dfrac{3\mathrm g}2\left(1-\dfrac{\sqrt3}2\right)$ m/s$^2$.}
\end{qoptions}
\qn{7} A rod is fixed between a vertical wall and a horizontal surface. A smooth ring of mass 1 kg is released from rest which can move along the rod as shown. At the release point spring is vertical and relaxed. The natural length of the spring is $\left(\sqrt3+1\right)$m (g = 10m/s$^2$) Rod makes an angle of 30$^\circ$ with the horizontal. Ring again comes to rest when spring makes an angle of 30$^\circ$ with the vertical\par
\qfig{p1_phy_s0_q07_1}
\begin{qoptions}
\opt{A}{Force constant of the spring is $\dfrac52\left(\sqrt3+1\right)$N / m}
\opt{B}{Maximum displacement of ring is $\dfrac2{\sqrt3-1}$m}
\opt{C}{Maximum extension in the spring is $\sqrt3-1$m}
\opt{D}{Normal reaction on ring due to rod when it again comes to rest is $\dfrac52\left(\sqrt3-1\right)$N}
\end{qoptions}
\qn{8} An elastic and light cord of length $l_0$ = 5.0 m and stiffness k = 50 n/mm is inside a fixed box. A small portion of length x$_0$ = 20 cm is left out of a hole made one of the walls. The cord fits tightly in the hole therefore when pulled frictional forces act on the cord. Total magnitude of these frictional forces is F = 25 N. Which of the following can not be work done in slowly pulling the complete length of the cord out of the box ?\par
\qfig{p1_phy_s0_q08_1}
\begin{qoptions}
\opt{A}{F = ($l_0-$x$_0$) + $\dfrac{\mathrm F^2}{\mathrm k}\left(1-\dfrac{\mathrm x_0}{2l_0}\right)$}
\opt{B}{F = ($l_0-$x$_0$) + $\dfrac{\mathrm F^2}{\mathrm k}\left(1-\dfrac{2\mathrm x_0}{l_0}\right)$}
\opt{C}{F = ($l_0-$x$_0$) + $\dfrac{2\mathrm F^2}{\mathrm k}\left(1-\dfrac{2\mathrm x_0}{l_0}\right)$}
\opt{D}{F = ($l_0-$x$_0$) $-$ $\dfrac{2\mathrm F^2}{\mathrm k}\left(1+\dfrac{2\mathrm x_0}{l_0}\right)$}
\end{qoptions}
\qn{9} A bead slides on a fixed frictionless wire bent into a horizontal semicircle of radius R$_0$ as shown in figure. In addition to any normal forces exerted by the wire, the bead is subjected to an external force that points directly away from origin and depends on distance r from the origin according to the formula $\vec{\mathrm F}=\mathrm F_0\left(\dfrac{\mathrm r}{\mathrm R_0}\right)^2\hat{\mathrm r}$ then which of the following are correct :\par
\qfig{p1_phy_s0_q09_1}
\begin{qoptions}
\opt{A}{Given force is a central force}
\opt{B}{Given force is a conservative force}
\opt{C}{Work done by external foce as bead leaves the track (starting from origin) is $\dfrac{8\mathrm F_0\mathrm R_0}3$}
\opt{D}{Speed of bead as it leaves the wire is $\sqrt{\mathrm v_0^2+\dfrac{18\mathrm F_0\mathrm R_0}{3\mathrm m}}$}
\end{qoptions}
\qn{10} One end of an light elastic cord is attached to a block of mass m and the other end is held keeping the cord relaxed and vertical. Relaxed length fo the cord is $\ell$ and there is a mark P at a height 0.8 $\ell$ from the lower end. If the upper end is slowly raised, the block leaves the ground, when the mark P reaches the point where the upper end of the relaxed cord initially was work done by the force applied by the hand will not be :\par
\qfig{p1_phy_s0_q10_1}
\begin{qoptions}
\opt{A}{$\dfrac{\mathrm{mg}\ell}{16}$}
\opt{B}{$\dfrac{\mathrm{mg}\ell}8$}
\opt{C}{$\dfrac{\mathrm{mg}\ell}5$}
\opt{D}{$\dfrac{\mathrm{mg}\ell}4$}
\end{qoptions}
\qn{11} A thin smooth massless bar is being made to spin with a contant angular velocity $\omega$ = 2 rad/s in a horizontal plane. A bead of mass 2 kg is sliding along the rod as shown. When the bead is at a distance r = 1 m from the axis of rotation of bar, its velocity relative to bar is 0.75 m/s. What is the acceleration of the bead with respect to ground at this moment (in m/s$^2$)\par
\qfig{p1_phy_s0_q11_1}
\qn{12} A particle A is moving in xy plane with the constant speed of 2$\pi$ m/s along the path x$^2$ + y$^2$ + 120y = 0. At time t = 0. When A is at the origin, another particle B starts moving from origin with constant acceleration is such a way that at time t = 5s velocities of both the particles are found to be equal. If sense of motion of A is clockwise, the average speed of B over first five second is k$\pi$. Find k.
\qn{13} A wedge is placed on a smooth horizontal surface. Wedge contains a circular quadrant of radius 25 cm as shown. An insect crawls on the circular part with a constant speed 1/2 m/s. A force F is applied on the wedge, so that it does not move. Find the value of F in newton when radial line of position of insect makes an angle 60$^\circ$ with horizontal.\par
\qfig{p1_phy_s0_q13_1}
\qn{14} In an XY horizontal plane a force field $\vec{\mathrm F}=-(40\mathrm N/\mathrm m)\left(\mathrm y\hat{\mathrm i}+\mathrm x\hat{\mathrm j}\right)$ is present where x and y are the coordinates of any point on the plane. A smooth rod AB is fixed in the plane as shown in the figure. A particle of mass 5 kg is to be released with a velocity in this force field such that it reaches to point B. Find the minimum velocity that must be imparted along the rod at A such that it reaches to B.\par
\qfig{p1_phy_s0_q14_1}
\qn{15} The block has mass M and rests on a surface for which the coefficient of friction $\mu$. If a force F = kt$^2$ is applied to the cable (see figure), the power developed by the force at t = t$_2$ is nearly \underline{\hspace{1cm}} kilowatt.\par
(Given : M = 20 kg, $\mu$ = 0.4, k = 40 N/s$^2$, t$_2$ = 3 sec.)\par
\qfig{p1_phy_s0_q15_1}
\qn{16} In the figure shown, the mass of the hanging blocks is m, while that of the block resting on the floor is 3m. The floor is horizontal and frictionless and all pulleys ideal. The system is initially held stationary, with the inclined thread making an angle $\theta$ = 30$^\circ$ with the horizontal. The blocks are now released from rest and allowed to move. The hanging blocks falls through a height (49/5) m before hitting the floor. It is found that the value of $\theta$ becomes 60$^\circ$, when the hanging block hits the floor. Find the speed with which the hanging block hits the floor.\par
\qfig{p1_phy_s0_q16_1}
\qn{17} A ball of mass m is suspended from the ceiling with the help of an elastic cord. An additional downward force applied on the ball, if increased gradually to a value F$_0$, the cord would break. The additional minimum constant force F that will break the cord is $\dfrac{\mathrm F_0}{\mathrm x}+(\mathrm{mg})\mathrm y$ . Then find x + y. (cord is massless and force is applied from equilibrium)
\qn{18} A massive homogeneous rope of length $\ell$ is suspended between two nails A and B driven in the same horizontal level. The rope makes angle $\theta$ with the horizontal at the nails and its lowest point is at a depth d below the nails. A bead threaded in the rope released from the nail A slides down the rope without friction. If the bead is so light that the shape of the rope remains unaffected, acceleration of the bead when it passes the lowest point of the rope is $\dfrac{\mathrm{x\,gd\,tan}\theta}{\ell}$. Then find x.\par
\qfig{p1_phy_s0_q18_1}
\qn{19} A block of mass 1 kg is pulled along the curve path ACB by a tangential force as shown in figure. The work done by the frictional force when the block moves from A to B is $-$5x Joule then find x. (Block is pulled very slowly)\par
\qfig{p1_phy_s0_q19_1}
\qn{20} At an instant t = 0, Block A of mass 1 kg is moving with speed 8 m/s towards right on rough surface of block B of mass 3 kg. Block B, which is placed on smooth horizontal surface is moving with speed 4 m/s towards right at same instant (t = 0). The net work done by the frictional force in long time is $-$x Joule then find x. (block A will not fall from block B)\par
\qfig{p1_phy_s0_q20_1}
\raggedcolumns
\qfullbegin
\papersec{1}
\partsec{2}{CHEMISTRY}
\qn{21} Select incorrect graphs, which are plotted between ionisation energy of 2nd period elements and their atomic number.\par
\qfig{p1_chem_s0_q21_1}
\qn{22} A molecule having formula PQ$_2$H$_8$ has following structural features (H is hydrogen atom)\par
(I) P and Q both belong to IInd period of periodic table.\par
(II) P and Q has two and three valence electrons respectively.\par
(III) There are four (2c---2e$^-$) bonds and four (3c---2e$^-$) bonds in molecule.\par
(IV) In P and Q both form atleast one (2c---2e$^-$) bond.\par
Based on above structural information which one of the following structure is not most appropriate for PQ$_2$H$_8$:\par
\qfig{p1_chem_s0_q22_1}
\qfullend
\qn{23} Choose the incorrect ionisation energy order.
\begin{qoptions}
\opt{A}{Li $<$ B $<$ Be $<$ C $<$ N $<$ O $<$ F}
\opt{B}{Na $<$ Mg $<$ Al $<$ Si $<$ S $<$ P $<$ Cl}
\opt{C}{K $<$ Ga $<$ Ca $<$ Ge $<$ As $<$ Se $<$ Br}
\opt{D}{Rb $<$ Sr $<$ In $<$ Sn $<$ Sb $<$ Te $<$ I}
\end{qoptions}
\qn{24} The following graph is given, between total energy and distance between the two nuclei for species H$_2^+$, H$_2$, He$_2^+$ \& He$_2$, which of the following statements is/are \textbf{incorrect:}\par
\qfig{p1_chem_s0_q24_1}
\begin{qoptions}
\opt{A}{He$_2^+$ is more stable than H$_2^+$.}
\opt{B}{Bond dissociation energy of H$_2^+$ is more than bond dissociation energy of He$_2^+$.}
\opt{C}{Since bond orders of He$_2^+$ and H$_2^+$ are equal hence both will have equal bond dissociation energy.}
\opt{D}{Bond length of H$_2^+$ is less than bond length of H$_2$ .}
\end{qoptions}
\qn{25} In which of the following arrangement, the order is according to the property indicated against it:
\begin{qoptions}
\opt{A}{I $<$ Br $<$ F $<$ Cl : Increasing EA}
\opt{B}{Li $<$ Na $<$ K $<$ Rb : Increasing metallic radius}
\opt{C}{B $<$ C $<$ N $<$ O : Increasing frist ionization energy}
\opt{D}{Al$^{+3}$ $<$ Mg$^{+2}$ $<$ Na$^+$ $<$ F$^-$ : Increasing ionic size}
\end{qoptions}
\qn{26} Which is/are \textbf{\underline{correct}} statement :
\begin{qoptions}
\opt{A}{In CF$_2$ = C = CF$_2$ molecule all the four fluorine atoms are not in the same plane.}
\opt{B}{Ka$_2$ fumaric acid is more than Ka$_2$ of maleic acid due to intra molecular hydrogen bonding in maleic acid.}
\opt{C}{The O $-$ O bond length in O$_2$[AsF$_4$] is shorter then KO$_2$.}
\opt{D}{The bond angle order in halogen $-$ S $-$ halogen is OSF$_2$ $<$ OSCl$_2$ $<$ OSBr$_2$}
\end{qoptions}
\qn{27} Following is the graph representing Electronegativity of certain elements A$_1$, A$_2$, A$_3$ , A$_4$, A$_5$. Which of the following option(s) representing corresponding elements can be \textbf{correct}?\par
\qfig{p1_chem_s0_q27_1}
\begin{qoptions}
\opt{A}{A$_1$ = \textbf{N}, A$_2$ = \textbf{P}, A$_3$ = \textbf{F}, A$_4$ = \textbf{Cl}, A$_5$ = \textbf{Al}}
\opt{B}{A$_1$ = \textbf{Cl}, A$_2$ = \textbf{P}, A$_3$ = \textbf{O}, A$_4$ = \textbf{N}, A$_5$ = \textbf{Al}}
\opt{C}{A$_1$ = \textbf{N}, A$_2$ = \textbf{P}, A$_3$ = \textbf{O}, A$_4$ = \textbf{Si}, A$_5$ = \textbf{Al}}
\opt{D}{A$_1$ = \textbf{N}, A$_2$ = \textbf{P}, A$_3$ = \textbf{F}, A$_4$ = \textbf{Al}, A$_5$ = \textbf{Sr}}
\end{qoptions}
\qn{28} Which of the following statmements is/are correct?
\begin{qoptions}
\opt{A}{In ClF$_3$ , the axial Cl $-$ F bond length is larger than equatorial Cl $-$ F bond length.}
\opt{B}{In SF$_4$, F $-$ S $-$ F equatorial angle is not 120$^\circ$ but 104$^\circ$ due to lp $-$ bp repulsions.}
\opt{C}{In [ICl$_4$]$^-$, Cl-I- Cl bond angle is 90$^\circ$.}
\opt{D}{In OBr$_2$, the bond angle is less than OCl$_2$.}
\end{qoptions}
\qn{29} Which of the following statements is/are not correct.
\begin{qoptions}
\opt{A}{Na$_2$O $<$ MgO $<$ ZnO $<$ P$_4$O$_{10}$ : Acidic property.}
\opt{B}{Na $<$ Si $>$ Mg $<$ Al : First ionisation potential.}
\opt{C}{F $>$ Cl $>$ Br : Electron affinity.}
\opt{D}{H$^-$ $>$ F$^-$ : Size}
\end{qoptions}
\qn{30} Select \textbf{CORRECT} statement(s).
\begin{qoptions}
\opt{A}{Correct boiling point order : H$_2$O $>$ HF $>$ NH$_3$}
\opt{B}{H$_2$ \& H$_3^+$ have same number of electrons}
\opt{C}{H$_3$O$^+$ \& NH$_4^+$ have same number of 2c-2e$^-$ bonds}
\opt{D}{NaHCO$_3$(s) \& KHCO$_3$(s) have same arrangement of hydrogen bonded HCO$_3^-$ ions.}
\end{qoptions}
\qn{31} Assuming all non--hydrogenic atoms undergo hybridization \& following are various categories of orbital overlapping.\par
sp$^3$--sp$^3$ axial ovelap (s) $\longrightarrow$ Type -- A\par
sp$^2$--sp$^2$ axial overlap (s) $\longrightarrow$ Type -- B\par
sp$^3$--sp$^2$ axial ovelap (s) $\longrightarrow$ Type -- C\par
p -- p lateral ovelap (p) $\longrightarrow$ Type -- D\par
p -- d lateral overlap (p) $\longrightarrow$ Type -- E\par
How many species which contain only 2 of these types (from A to E) of overlapping of orbitals.\par
(i) H$_2$CO$_3$ \hfill (ii) B$_3$N$_3$H$_3$F$_3$ \hfill (iii) SNF$_3$ \hfill (iv) H$_2$S$_4$O$_6$
\qn{32} Calculate the 10$^\mathrm{th}$ ionisation energy of Neon in Mega Joules/mole.\par
[Use N$_\mathrm A$ = 6 $\times$ 10$^{23}$ \& 1 eV = 1.6 $\times$ 10$^{-19}$ J]
\qn{33} Find number of compound in which both underlined atoms are \textbf{NOT} directly bonded but connected via bridging atom(s).\par
$\underline{\mathrm C}_2\underline{\mathrm{Cl}}_6$, $\underline{\mathrm{Al}}_2\underline{\mathrm{Cl}}_6$, $\underline{\mathrm{Si}}_2\underline{\mathrm{Cl}}_6$, $\underline{\mathrm{Be}}_2\underline{\mathrm{Cl}}_4$, $\underline{\mathrm I}_2\underline{\mathrm{Cl}}_6$, $\underline{\mathrm B}_2\underline{\mathrm H}_6$, H$_2$S$_2$O$_6$, Na$_2$[$\underline{\mathrm{Pt}}_2\underline{\mathrm{Cl}}_6$], ($\underline{\mathrm{Be}}\underline{\mathrm H}_2$)$_\mathrm n$
\qn{34} The outermost \& the penultimate shell contain $\dfrac14$th \& $\dfrac23$rd of their maximum capacity of electrons in an element. Calculate the effective nuclear charge for the valence shell electron as per Slater's rule.
\qn{35} Few processes are as follow :\par
(i) O$_2\longrightarrow$ O$_2^+$ \hfill (ii) N$_2^-\longrightarrow$ N$_2$ \hfill (iii) PCl$_5$ + Cl$^-\longrightarrow$ PCl$_6^-$\par
(iv) N$_2$O$_4$(s)$\longrightarrow$ P (Cation) + Q (Anion)\par
If number of unpaired electron change in (i) is (a) where a $>$ 0\par
Change in bond order in (ii) is (b)\par
Increase in number of 90$^\circ$ bond angles in (iii) is (c)\par
Bond order of N--O in (Q) is (d) \& Bond order of N--O in (P) is (e). Then specify 5a + 2b + c + 6d + 2e.
\qn{36} When electrons (one from each atom) are removed from 0.5 mol of A and given them (one to each atom) to 0.5 mol of B, 40 kJ of energy is absorbed. When electrons (one from each atom) are removed from 0.5 mol of B and given them (one to each atom) to 0.5 mol of A, 52 kJ energy is released. `When electrons (two from each atom) are removed from 0.5 mol of A and given them (one to each atom) to one mole of B, 280 kJ energy is absorbed\par
IE$_1$ of A = 400 KJ/mol\par
Electron gain enthalpy of A = $-$352 KJ/mol\par
Sample of A and B are taken in gas phase.\par
(i) Ionisation energy of B (in KJ/mol) = `X'\par
(ii) Second ionisation energy of A (in KJ/mol) = ``Y''\par
(iii) Electron gain enthalpy of B (in KJ/mol) = ``Z''\par
\textbf{Then find the value of (X + Y + Z).}
\qfullbegin
\qn{37} The following question refer to the twelve compounds given below.\par
\qfig{p1_chem_s0_q37_1}
\qfullsplit
How many compounds may serve only as H-bond acceptors ?
\qfullend
\qn{38} An ionic compound AX$_3$ (s) is formed between a metal A and a non-metal X (outermost shell configuration of X = ns$^2$ np$^5$). Find the magnitude of enthalpy of formation (Kcal/mole) of AX$_3$(s) with the help of the following data.\par
(Non-metal X is found to exist in nature as a diatomic gas)\par
$\Delta$H$_\mathrm{Sublimation}$ A(s) = 100 Kcal/mol\par
$\Delta$H$_{\mathrm{I.E}_1}$ A(g) = 60 Kcal/mol\par
$\Delta$H$_{\mathrm{I.E}_2}$ A (g) = 150 Kcal/mol\par
$\Delta$H$_{\mathrm{I.E}_3}$ A(g) = 280 Kcal/mol\par
$\Delta$H$_\mathrm{diss}$ X$_2$(g) = 80 Kcal/mol\par
$\Delta$H$_\mathrm{e.g}$ X(g) = $-$ 110 Kcal/mol\par
$\Delta$H$_\mathrm{Lattice\;energy}$ A X$_3$(s) = $-$ 470 Kcal/mol
\qn{39} Calculate the total number of valence electrons, lone pairs and p$\pi$ -- p$\pi$ bonds. In the given structure if Y = P atom , X = N atom, Z = Cl atom\par
\qfig{p1_chem_s0_q39_1}
\qn{40} IE$_1$ (kJ/mol) of 2nd period elements are mentioned below with some blank/vacant data spaces.\par
{\footnotesize\setlength{\tabcolsep}{2pt}\noindent\begin{tabular}{cccccccc}
Li & Be & B & C & N & O & F & Ne \\
520 & \underline{``P''} & 800 & \underline{``Q''} & 1402 & 1314 & \underline{``R''} & 2080 \\
\end{tabular}}\par
Select correct number of statement(s) regarding above data table. (Consider all processes are gas phased)\par
(i) $|$EA$|$ of N$^+$ is more than 1402 kJ/mol\par
(ii) Value of ``R'' is more than 1402 kJ/mol\par
(iii) Value of P is more than $|$EA$|$ of C$^+$\par
(iv) Q $>$ 800 kJ/mol\par
(v) $\dfrac12$ mole O$^+\to$ O released energy = S ; $\dfrac12$ mole Ne$^+\to$ Ne released energy = T, then ($|$S$|$ + $|$T$|$ $>$ 2R) is correct.)
\raggedcolumns
\papersec{1}
\partsec{3}{MATHEMATICS}
\qn{41} Three equal circles of radius 1 touch one another. Radius of the circle touching all the three circles can be :
\begin{qoptions}
\opt{A}{$\dfrac{2-\sqrt3}{\sqrt3}$}
\opt{B}{$\dfrac{\sqrt3-\sqrt2}{\sqrt2}$}
\opt{C}{$\dfrac{2+\sqrt3}{\sqrt3}$}
\opt{D}{$\dfrac{\sqrt3+\sqrt2}{\sqrt2}$}
\end{qoptions}
\qn{42} If 2 cos $\theta$ + sin $\theta$ = 1, then the value of 4 cos $\theta$ + 3 sin $\theta$ can be equal to :
\begin{qoptions}
\opt{A}{3}
\opt{B}{$-$ 5}
\opt{C}{$\dfrac75$}
\opt{D}{$-$ 4}
\end{qoptions}
\qn{43} The sides of a $\Delta$ABC satisfy the equation, 2a$^2$ + 4b$^2$ + c$^2$ = 4ab + 2ac. Then :
\begin{qoptions}
\opt{A}{the triangle is isosceles}
\opt{B}{the triangle is obtuse}
\opt{C}{B = cos$^{-1}\dfrac78$}
\opt{D}{A = cos$^{-1}\dfrac14$}
\end{qoptions}
\qn{44} If in a $\Delta$ABC, cosA . cosB + sinA sinB sin2C = 1 then,
\begin{qoptions}
\opt{A}{$\Delta$ABC is isosceles but not right angled}
\opt{B}{$\Delta$ABC is acute angled}
\opt{C}{$\Delta$ABC is right angled}
\opt{D}{least angle of the triangle is $\dfrac\pi4$}
\end{qoptions}
\qn{45} $\displaystyle\prod_{\mathrm r=1}^{11}\sec\left(\dfrac{\mathrm r\pi}{11}\right)$ is divisible by :
\begin{qoptions}
\opt{A}{4}
\opt{B}{16}
\opt{C}{64}
\opt{D}{256}
\end{qoptions}
\qn{46} The value of $\dfrac{3+\cot76^\circ\cot16^\circ}{\cot76^\circ+\cot16^\circ}$ is :
\begin{qoptions}
\opt{A}{cot 44$^\circ$}
\opt{B}{tan 46$^\circ$}
\opt{C}{tan 44$^\circ$}
\opt{D}{cot 46$^\circ$}
\end{qoptions}
\qn{47} The sines of two angles of a triangle are equal to $\dfrac5{13}$ and $\dfrac{99}{101}$ . The cosine of the third angle can be :
\begin{qoptions}
\opt{A}{$\dfrac{245}{1313}$}
\opt{B}{$\dfrac{255}{1313}$}
\opt{C}{$\dfrac{735}{1313}$}
\opt{D}{$\dfrac{765}{1313}$}
\end{qoptions}
\qn{48} In a $\Delta$ABC, if the median, bisector and altitude drawn from the vertex A divide the angle at the vertex into four equal parts then the angle of the $\Delta$ABC can be equal to :
\begin{qoptions}
\opt{A}{$\dfrac\pi8$}
\opt{B}{$\dfrac\pi2$}
\opt{C}{$\dfrac{3\pi}8$}
\opt{D}{$\dfrac\pi5$}
\end{qoptions}
\qn{49} Sum of all the solutions of the equation cos 4x + 6 = 7 cos 2x in [0, 314] is k$\pi$. then :
\begin{qoptions}
\opt{A}{sum of digits of k is 18}
\opt{B}{k is divisible by 25}
\opt{C}{sum of digits of k is 15}
\opt{D}{k is divisible by 7}
\end{qoptions}
\qn{50} If N represents the number of real solutions of the equation sin x + sin y = sin(x+y) and $|$x$|$ + $|$y$|$ =1 simultaneously, then N is greater than
\begin{qoptions}
\opt{A}{2}
\opt{B}{4}
\opt{C}{6}
\opt{D}{8}
\end{qoptions}
\qn{51} In an isosceles $\Delta$ABC if the altitudes intersect on the inscribed circle then the secant of the vertical angle `A' is :
\qn{52} The value of $\cos\dfrac\pi{20}.\cos\dfrac{3\pi}{20}.\cos\dfrac{7\pi}{20}.\cos\dfrac{9\pi}{20}+\cos\dfrac\pi{15}.\cos\dfrac{2\pi}{15}.\cos\dfrac{4\pi}{15}.\cos\dfrac{8\pi}{15}$ is
\qn{53} The value of $\sec\dfrac\pi7\left(\dfrac{3-\tan^2\dfrac\pi7}{1-\tan^2\dfrac\pi7}\right)$ is
\qn{54} If the circum centre of the $\Delta$ABC lies on its incircle then the value of [R/r] is where [ . ] is GI function :
\qn{55} Let a, b and c be the sides of a $\Delta$ABC. If a$^2$, b$^2$ and c$^2$ are the roots of the equation x$^3-$8x$^2$+6x $-$ 4 = 0, then the value of $\dfrac{\cos\mathrm A}{\mathrm a}+\dfrac{\cos\mathrm B}{\mathrm b}+\dfrac{\cos\mathrm C}{\mathrm c}$ is equal to :
\qn{56} The exact value of $\dfrac{96\sin80^\circ\sin65^\circ\sin35^\circ}{\sin20^\circ+\sin50^\circ+\sin110^\circ}$ is equal to :
\qn{57} The value of $\cos\dfrac{2\pi}{28}\mathrm{cosec}\dfrac{3\pi}{28}+\cos\dfrac{6\pi}{28}\mathrm{cosec}\dfrac{9\pi}{28}+\cos\dfrac{18\pi}{28}\mathrm{cosec}\dfrac{27\pi}{28}$ is
\qn{58} Number of real roots of the equation sec$\theta$ + cosec$\theta$ = $\sqrt{15}$ lying between 0 and $\pi$ is :
\qn{59} The number of real solutions of the equation ($\sin\theta$+2)($\sin\theta$+3)($\sin\theta$+4)=6 in $\left[-\pi,\dfrac{5\pi}2\right)$ is :
\qn{60} If A = $\sin\dfrac{2\pi}7+\sin\dfrac{4\pi}7+\sin\dfrac{8\pi}7$ and B = $\cos\dfrac{2\pi}7+\cos\dfrac{4\pi}7+\cos\dfrac{8\pi}7$ , then the value of A$^2$ + B$^2$ is :
\end{multicols}
\end{document}
